\documentclass[11pt,a4paper]{article}
\usepackage[margin=1in]{geometry}
\usepackage{hyperref}
\usepackage{enumitem}
\usepackage{graphicx}
\usepackage{xcolor}
\usepackage{titling}

\title{Smart Library Seat Management System}

\author{
  Ayush Singh Sikarwar \and Satvik Ganda \and Aryan Bansal \and Rishit Parnami
}

\date{\today}

\begin{document}
\maketitle
\begin{center}
\large \textbf{Project Report: Final Stage} \\
\vspace{0.5cm}
UCS503: Software Engineering \\
Thapar Institute of Engineering and Technology
\end{center}

\section{Executive Summary}
The Smart Library Seat Management System is a comprehensive web-based platform designed to optimize physical study space allocation. By introducing real-time occupancy tracking and automated policy enforcement, the system drastically reduces seat-hunting friction and ensures fair usage of library resources.

\section{System Architecture}
The system employs a scalable three-tier architecture:
\begin{itemize}
    \item \textbf{Frontend:} A responsive single-page application enabling students to interact with the library floor map and manage their session statuses.
    \item \textbf{Backend API:} A RESTful service handling authentication, session timers, and seat allocation rules.
    \item \textbf{Database:} A relational database maintaining persistent state for users, physical layout maps, and historical analytics.
\end{itemize}

\section{Key Features Delivered}
\begin{enumerate}
    \item \textbf{Interactive Floor Map:} Real-time visualization of seat availability across library zones.
    \item \textbf{Temporary Away Permit:} Allows students to temporarily leave their seat (e.g., 30 mins limit), releasing it automatically if they overstay.
    \item \textbf{Smart Seat Finder:} AI-driven recommendations based on student preferences (e.g., power outlets, quiet zones).
    \item \textbf{Admin Dashboard:} Provides library administrators with occupancy metrics, allowing them to manage policies and monitor daily usage trends.
\end{enumerate}

\section{Evaluation \& Results}
The pilot testing revealed a significant reduction in the average time students spent searching for seats. Furthermore, the automated release mechanism for expired 'Temporary Away' sessions recovered a measurable percentage of seating capacity that would otherwise remain blocked by unattended belongings.

\section{Conclusion \& Future Work}
The deployed system meets all core objectives of improving transparency and fairness in library seating. Future iterations could integrate physical IoT sensors for automated occupancy validation and construct a 3D digital twin of the library environment for enhanced spatial visualization.

\end{document}
